\section{The Claim That Crosses the Boundary}
\label{sec:ch15-the-claim-that-crosses}
\index{header forwarding!the router's \code{headers} block|(}%

Chapter~\ref{ch:how-a-federated-query-actually-runs} put a logger inside two
subgraphs and found that the router forwards no client header at all. That was a
measurement in search of a consequence for eight chapters. Here is the
consequence: two of Mosaic's three enforcement layers run inside a subgraph, and
a subgraph with no header has no caller.

The smallest thing that changes it is five lines:

\begin{minted}[fontsize=\footnotesize,breaklines]{yaml}
headers:
  all:
    request:
      - op: propagate
        named: Authorization
\end{minted}

\code{propagate} and \code{set} are the two operations the block
has~\autocite{cosmodocs2026headers}. \code{propagate} moves a header the client
sent; \code{set} writes one the router computes.

\subsection*{Which of the two arrangements, and why}

\index{header forwarding!the token or a claim derived from it}%
WunderGraph documents the other arrangement and recommends
it~\autocite{cosmodocs2026headers}: do not forward the token, and instead set a
header derived from a claim.

\begin{minted}[fontsize=\footnotesize,breaklines]{yaml}
      - op: set
        name: "X-Mosaic-Customer"
        expression: "request.auth.isAuthenticated ? request.auth.claims.sub : ''"
\end{minted}

It is the better arrangement on every axis but one. It is smaller on the wire.
It keeps the token in one process, so a subgraph log that dumps headers is not a
credential leak. It means a subgraph needs no JWT library, no signing key and no
clock. And every subgraph gets a plain string it can read without ceremony.

The axis it loses on is the one that matters here: it requires the subgraph to
trust a header, which is only safe if nothing but the router can send one.
Mosaic publishes all seven services on host ports. A header that says
\code{X-Mosaic-Customer} is a header anybody with \code{curl} can also say, and
the whole model would come apart on a request the router never saw.

So Mosaic forwards the token itself and every service validates it again. The
cost is a signature verification per subgraph per request, which is real and
small. The cost of the other choice is that the security of the graph depends on
a network boundary that nothing in this repository creates or checks, and which
somebody re-opens the first time they add a deployment target. I would take the
documented arrangement in a heartbeat behind a mesh I controlled. I am not going
to take it behind seven published ports and a comment.

\index{header forwarding!one rule covers both transports}%
One detail worth knowing before you need it, measured on a two-service probe
rather than on Mosaic: that single \code{propagate} entry covers the WebSocket
upgrade as well as an ordinary HTTP fetch. There is no second,
subscription-specific header block for it. Mosaic cannot show you that, because
by the end of this chapter no subscription it serves carries a token at all,
which is section~\ref{sec:ch15-a-connection-with-no-identity}'s subject.

\subsection*{What a subgraph does with it}

\index{Mosaic.ServiceDefaults@\code{Mosaic.ServiceDefaults}!the security defaults}%
Validating a bearer token is the same work in every service and nothing any two
of them have to agree about, so by chapter~\ref{ch:strangling-the-monolith}'s
test it is platform and lives in the shared library. Four of the seven call it.

Two lines in there are worth reading twice. The first:

\begin{minted}[fontsize=\footnotesize,breaklines]{csharp}
options.MapInboundClaims = false;
\end{minted}

Left at its default of true, ASP.NET Core rewrites well-known JWT claim names
into the long-form URIs of the old WS-Federation claim types, so the \code{sub}
the token carries arrives under a name ending in \code{nameidentifier} and a
resolver looking for \code{sub} finds nothing. The token is correct, the
signature is valid, the request is authenticated, and the claim is missing under
the name it was written with. Turning the map off makes the claim names in a
resolver the claim names in the token, which are also the names the router
reads.

The second is a policy that does not work:

\begin{minted}[fontsize=\footnotesize,breaklines]{csharp}
policy.RequireClaim("scope", "orders:read")   // wrong
\end{minted}

A scope claim is one string holding every granted scope separated by spaces, so
that comparison puts \code{"orders:read reviews:write"} up against
\code{"orders:read"} and refuses a token that plainly holds the scope. The
failure is a \code{403} on a correct token, which is a bad half hour. Mosaic's
policies split on the space, once, in a shared helper, because the alternative
is getting it right in three services.

\subsection*{Two more middleware in the pipeline}

\index{request pipeline!what authorization adds}%
Turning authorization on changes something
chapter~\ref{ch:the-life-of-a-request} counted. That chapter walked thirteen
request middleware and both verification scripts have asserted the list in order
ever since. A service that calls \code{AddAuthorization()} assembles fifteen,
and the two new ones do not go on the end:

\begin{minted}[fontsize=\footnotesize,breaklines]{text}
 5. DocumentParserMiddleware
 6. HotChocolate.Authorization.Pipeline.PrepareAuthorization
 7. DocumentValidationMiddleware
 8. HotChocolate.Authorization.Pipeline.AuthorizeRequest
 9. CostAnalyzerMiddleware
\end{minted}

One in front of validation and one behind it. That placement is why the
\code{[Authorize]} attribute has an \code{apply} argument at all: set to
\code{Validation}, a refusal happens in the request pipeline before any resolver
runs, which is how a whole operation gets rejected rather than one field nulled.
Mosaic uses the default, \code{BeforeResolver}, everywhere, and the difference
is visible from outside. At the default a refused field is an error with a
\code{path} and an HTTP \code{200}. At \code{Validation} it is an HTTP
\code{401} with no \code{path} and no \code{data} key at all.

Chapter~\ref{ch:the-life-of-a-request}'s thirteen is still right for the service
that chapter measured, and four of the seven still assemble exactly thirteen.
The gate holds both lists now, per subgraph, and asserts the order of each. It
would have been easier to assert \enquote{thirteen or fifteen} and that is
exactly the kind of loosening this book's own rules forbid.

\medskip
A token at the router, a token at the subgraph, and a policy in front of a
resolver. That is two of the three layers working. The third one is the
interesting one, and the field that makes it necessary is the one this graph
grew two chapters ago to be helpful.

\index{header forwarding!the router's \code{headers} block|)}%
